\documentclass[11pt]{article}
\usepackage[margin=1.1in]{geometry}
\usepackage{amsmath, amssymb, amsthm}
\usepackage{mathtools}
\usepackage{booktabs}
\usepackage{enumitem}
\usepackage{hyperref}
\usepackage{parskip}
\usepackage{xcolor}
\hypersetup{colorlinks=true, urlcolor=blue, linkcolor=black}

\title{A Log-Linear IRT Accumulator Model\\ of Early Word Learning}
\author{Draft specification (revised: 2PL extension)}
\date{\today}

\newcommand{\E}{\mathbb{E}}
\newcommand{\Var}{\mathrm{Var}}
\newcommand{\Prob}{\mathbb{P}}
\newcommand{\Bern}{\mathrm{Bernoulli}}
\newcommand{\Norm}{\mathcal{N}}
\newcommand{\HN}{\mathrm{HalfNormal}}
\newcommand{\logit}{\mathrm{logit}}
\newcommand{\sigmoid}{\sigma}

\begin{document}
\maketitle

\section{Research questions driving the modeling}

The model is organized to answer three research questions. The baseline
is the minimal specification that lets each be posed and answered;
each variant in the nested family (§\ref{sec:nested}) sharpens one of
them. A fourth question — when does accumulation start? — is posed in
the appendix because the start-time parameter \(s\) is
near-non-identifiable in practice (see §\ref{sec:s-non-id}).

\begin{enumerate}[leftmargin=1.3em, label=\textbf{RQ\arabic*.}]
\item \textbf{Is there acceleration in the accumulation process?}
  McMurray (2007) and Mitchell \& McMurray (2009) argued from theory
  that the empirical ``vocabulary explosion'' requires acceleration in
  the threshold distribution \(g(\tau)\): pure constant-rate
  accumulation produces growth that's too smooth, and leveraged
  learning shifts the curve in time but does not create acceleration
  that isn't already present. Bergelson (2020) made the empirical
  case: older infants are ``better learners.'' We formalize this
  question via \(\delta\): in our linear predictor, time enters as
  \((1 + \delta)\,\log\!((t-s)/a_0)\), where \(\delta = 0\)
  corresponds to McMurray's pure linear accumulator (one log-token of
  evidence per log-age unit, no acceleration). Posterior \(\delta > 0\)
  with CrI excluding zero is a quantitative confirmation of the
  McMurray–Mitchell acceleration claim.

\item \textbf{How well does frequency reconstruct word difficulty?}
  Braginsky et al.\ (2019) showed frequency is only modestly related
  to AoA within lexical classes. We test this within our model by
  asking how much of word-level residual difficulty \(\psi_j\) is
  explained by \(\log p_j\), and whether the per-class slope
  \(\beta_c\) on \(\log p_j\) deviates from the unit-accumulator
  baseline \(\beta_c = 1\).

\item \textbf{How much of child-level outcome variance is input-driven?}
  This is the \(\sigma_\alpha\) question: once we pin the variance of
  CDS input rate \(r_i\) from external data, how much residual child
  variance must be absorbed by a learning-efficiency term
  \(\alpha_i\)? Posterior on
  \(\pi_\alpha = \sigma_\alpha^2 / (\sigma_r^2 + \sigma_\alpha^2)\)
  reports the efficiency share of intercept variance.
\end{enumerate}

Each RQ maps onto a specific posterior quantity, recovered as the
nested family is built up:

\begin{table}[h]
\centering
\small
\begin{tabular}{@{}lll@{}}
\toprule
RQ & Quantity & How answered \\
\midrule
1 & \(\delta\) & Posterior in M2 vs.\ M1; baseline (M1) pins \(\delta=0\). \\
2 & \(\beta_c\), \(r(\hat\psi_j, \log p_j)\) &
  Posterior in M4 (per-class \(\beta_c\)); residual correlation in M3. \\
3 & \(\pi_\alpha = \sigma_\alpha^2/(\sigma_r^2 + \sigma_\alpha^2)\) &
  Derived posterior in M3+; \(\sigma_r^2\) pinned externally. \\
\bottomrule
\end{tabular}
\caption{Mapping from research questions to model quantities and
nested-family stages (M1--M5).}
\label{tab:rq-map}
\end{table}

\section{Motivating the model: how much variation is there to explain?}
\label{sec:precursor}

Before fitting any accumulator model, it is worth asking how much
between-child variation in language ability actually exists in the
data, and how that variation relates to age-driven development. We
ran the simplest possible 1-PL IRT —
\(P(y_{aj} = 1) = \sigma(\theta_a - \beta_j)\), with one ability
\(\theta_a\) per admin (\emph{not} per child) and one difficulty
\(\beta_j\) per word — fit via \texttt{lme4::glmer} on the Wordbank
English longitudinal sample. No time term, no frequency, no per-child
structure. Just: how variable is ability across admins, and how does
the population mean change with age?

\paragraph{Two reference quantities.} The population-mean
\(\theta\) rises by \(\hat \beta_{\mathrm{age}} = 0.47\) logits/month
on this sample (a regression of admin-level \(\theta_a\) on age).
Pooled across all admins, \(\mathrm{SD}(\theta_a) = 3.25\) logits;
the residual within-age SD (after subtracting the linear age trend)
is 1.95 logits.

\paragraph{Translating SD into months of development.} A natural
unit for individual differences is ``how many months of typical
development does \(\pm 1\) SD across kids correspond to?'' On this
sample,
\[
\frac{\mathrm{SD}_{\mathrm{within-age}}(\theta)}
     {\hat \beta_{\mathrm{age}}}
= \frac{1.95}{0.47}
\approx 4.2~\text{months of development per SD}.
\]
A child one SD above the within-age mean looks, on the IRT scale,
like a child of average ability 4 months older.

\paragraph{Variation grows with development.} The within-age SD is
not constant: it rises from \(\sim 1.0\) logit at 12--14~mo
(\(\sim 2\) months equivalent) to \(\sim 2.7\) at 26--30~mo
(\(\sim 5.8\) months equivalent). Children \emph{fan out} rather
than converge as they develop — a finding that any model of
individual differences must accommodate.

\paragraph{What this implies for the model.} The variation we need
to explain is large and increasing with age. A pure accumulator
with constant per-child input rate cannot produce growing spread; we
need either age-varying efficiency (per-child slope deviation
\(\zeta_i\)) or some other mechanism. The remainder of the
explainer builds the nested model family that makes structural
choices about \emph{which} mechanism. Figures
\texttt{precursor\_trajectories.png} and
\texttt{precursor\_variation.png} (in
\texttt{model/figs/longitudinal/}) and the script
\texttt{model/scripts/precursor\_irt.R} reproduce these numbers.

\section{One-paragraph summary}

A child \(i\) of age \(a_i\) has, on average, heard word \(j\) approximately
\(T_{ij} = r_i\,p_j\,H\,(a_i - s)\) times, where \(r_i\) is the child's
tokens-per-hour of CDS, \(p_j\) is the word's corpus probability, \(H\) is
waking hours per month, and \(s\) is the global start time for
accumulation. The child produces word \(j\) on the CDI with probability
\(\sigmoid(\log T_{ij} + \log \alpha_i - \psi_j)\) under a canonical
Rasch/IRT logit link (\(\sigmoid(x) = 1/(1+e^{-x})\)), where \(\alpha_i\)
is the child's learning-efficiency multiplier and \(\psi_j\) is the
word's log-threshold. Every quantity lives on a log-token scale with
\(O(1)\) priors, so Stan is well-conditioned. The crucial move is that
\(r_i\)'s population distribution is pinned by external data
(Sperry/HR/WF), so \(\sigma_\alpha\) is identifiable and Question 4 has an
answer. Questions 1--3 are posterior summaries of the baseline;
planned variants sharpen them.

\section{Quantities}

\begin{center}
\small
\begin{tabular}{@{}lll@{}}
\toprule
Symbol & Meaning & Units \\
\midrule
\(r_i\)         & Child \(i\)'s CDS input rate                & tokens / hour \\
\(\alpha_i\)    & Child \(i\)'s learning-efficiency multiplier & dimensionless \\
\(p_j\)         & Word \(j\)'s corpus probability (CHILDES)   & dimensionless, \(\sum_j p_j = 1\) \\
\(H\)           & Waking hours per month (fixed)              & hours / month \\
\(a_i\)         & Age at CDI assessment                       & months \\
\(s\)           & Global start time (accumulation onset)      & months \\
\(\psi_j\)      & Log expected tokens of word \(j\) to acquisition & log-tokens \\
\(\delta\)      & Age rate-change exponent                    & dimensionless \\
\(\lambda_j\)   & Word \(j\)'s 2PL discrimination             & dimensionless ($>0$) \\
\(a_0\)         & Reference age (fixed)                       & months \\
\(y_{ij}\)      & Child \(i\) produces word \(j\) on CDI      & \(\{0, 1\}\) \\
\bottomrule
\end{tabular}
\end{center}

\textbf{Fixed constants.} \(H = 12 \times 30.44 \approx 365\) hours/month
(12 waking hours/day). \(a_0 = 20\) months.

\section{Generative model}

\subsection{Linear predictor}

For child \(i\), word \(j\), age \(a_i > s\):
\begin{equation}
\eta_{ij}
= \lambda_j \cdot \Big[ \log r_i + \log \alpha_i
+ \log p_j + \log H
+ (1+\delta)\,\log\!\left(\tfrac{a_i - s}{a_0}\right)
- \psi_j \Big].
\label{eq:eta}
\end{equation}

Here \(\lambda_j \ge 0\) is a word-specific \emph{discrimination} parameter
(2PL). With \(\lambda_j \equiv 1\) the model collapses to the Rasch form
\(\eta_{ij} = \log(r_i \alpha_i p_j H (a_i - s)) - \psi_j\) — the log
ratio of tokens heard to tokens needed, with unit logit-per-log-token
gain. With \(\lambda_j\) free, each word has its own gain; high
\(\lambda_j\) produces a sharp transition near \(\psi_j\), low
\(\lambda_j\) a gradual one. Section \ref{sec:twopl} motivates the
extension in terms of the accumulator and in terms of model fit.

\subsection{Likelihood}

Canonical Rasch/IRT logit link:
\begin{equation}
y_{ij} \mid \eta_{ij} \sim \Bern\big(\sigmoid(\eta_{ij})\big),
\qquad \sigmoid(x) = \frac{1}{1 + e^{-x}}.
\label{eq:likelihood}
\end{equation}

The logit link is the standard ``log-linear IRT'' choice and is
numerically well-behaved in Stan. \(\psi_j\) is interpretable as the
log-threshold in log-tokens on the logit scale (the scale factor is
absorbed implicitly; \(\exp(\psi_j)\) is proportional to, not exactly,
the expected number of tokens to 50\%-acquisition).

Edge case: when \(a_i \le s\), set \(\Prob(y_{ij}=1) = \varepsilon\) and
exclude these cells from the log-age term (in the code this is
implemented by flooring $a_i - s$ at $0.01$ inside the $\log$).

\subsection{Priors}

\paragraph{Child-level latent (collapsed).}
The likelihood depends on \(\log r_i\) and \(\log \alpha_i\) only through
their sum, so we sample a single child-level latent
\[ \xi_i \coloneqq \log r_i + \log \alpha_i \sim \Norm\!\big(\mu_r,\; \sigma_r^2 + \sigma_\alpha^2\big). \]
Here \(\mu_r = \mathrm{mean}(\log r)_{\text{Sperry pool}} \approx 7.34\)
(geometric mean \(\approx 1537\) tokens/hour) and
\(\sigma_r^2 = \Var(\log r)_{\text{Sperry pool}} \approx 0.534^2\) are
\emph{pinned} from external data (this is the identifiability lever
for RQ4). The pool is the \(n=42\) dyads with non-NA
\(\texttt{adult\_child\_tokens\_hr}\) in
\texttt{data/raw\_data/sperry/hourly\_tokens\_Sperry\_HartRisley.csv}; in
practice these are all from Sperry et al.\ (2019), since Hart \& Risley
(1995) and Weisleder \& Fernald (2013) reported only mother-only and
all-adult counts, respectively. See \texttt{input\_estimation/} for the
full validation set and a sensitivity-range justification.
\(\sigma_\alpha\) is free:
\begin{equation}
\sigma_\alpha \sim \HN(0, 1).
\end{equation}

The per-child \(\log \alpha_i\) is a \emph{derived} quantity with closed-form
posterior conditional on \(\xi_i\):
\[
\log \alpha_i \mid \xi_i \;\sim\;
\Norm\!\left(\frac{\sigma_\alpha^2}{\sigma_r^2 + \sigma_\alpha^2}\,(\xi_i - \mu_r),
\;\frac{\sigma_r^2 \sigma_\alpha^2}{\sigma_r^2 + \sigma_\alpha^2}\right),
\]
reported in generated quantities. (Rationale for this collapsed form is
in \S\ref{sec:identifiability}: sampling both latents individually
over-parameterizes the model and creates bad posterior geometry in HMC.)

\paragraph{Word thresholds by lexical class.}
\begin{equation}
\psi_j \mid \mathrm{class}(j) = c \sim \Norm(\mu_c, \tau_c^2),
\quad \mu_c \sim \Norm(8, 3^2), \quad \tau_c \sim \HN(0, 1).
\end{equation}

\textbf{Note on \(\mu_c\).} The earlier value \(\log 1000\) was a
back-of-envelope centering. A more principled choice: tokens heard of a
typical word by age 20~mo spans roughly \(e^5\) (rare words) to \(e^{12}\)
(function words), so \(\mu_c \sim \Norm(8, 3^2)\) is weakly informative
across the relevant range and dominated by the data.

\paragraph{Global start time (RQ2).}
\begin{equation}
s \sim \Norm(4.5, 2^2) \text{ truncated to } [0, 15].
\end{equation}

One $s$ for the whole model, per your decision. The prior is centered at
4.5~mo — M\&P's comprehension estimate was \(\sim\)2~mo; production is
likely later. \(s\) close to 0 would indicate accumulation begins near
birth.

\paragraph{Age rate-change (RQ3).}
\begin{equation}
\delta \sim \Norm(0, 0.5^2).
\end{equation}

\paragraph{Word discrimination (2PL).}
\begin{equation}
\log \lambda_j \sim \Norm(0, \sigma_\lambda^2),
\quad \sigma_\lambda \sim \HN(0, 1).
\end{equation}

Setting \(\sigma_\lambda \to 0\) (via a near-zero prior SD passed as data)
pins every \(\lambda_j = 1\) and recovers the Rasch model. This is the
toggle used in code: one Stan file serves both Rasch and 2PL variants.
Scale identifiability (the classic \(\lambda_j \leftrightarrow 1/\lambda_j\),
\(\psi_j, \xi_i \leftrightarrow \lambda_j \psi_j, \lambda_j \xi_i\)
ambiguity of 2PL) is resolved by our pinning of \(\sigma_r\): the
\(\xi\) scale is fixed externally, so \(\psi_j\) and \(\lambda_j\) are
separately identified.

\(\delta = 0\) is the pure cumulative null. \(\delta > 0\) indicates older
children accumulate faster — a formal test of the Bergelson
comprehension-boost claim.

\section{How each RQ is answered}

\subsection*{RQ1 — is there acceleration?}

Single parameter \(\delta\), entering the time term as
\((1 + \delta)\,\log\!((t-s)/a_0)\).  \(\delta = 0\) is McMurray's
pure linear accumulator: each log-age unit contributes one log-token of
evidence. \(\delta > 0\) means evidence accumulates super-linearly in
log-age (or equivalently: tokens grow as \(t^{1+\delta}\) on the
linear scale, faster than the linear-accumulation baseline). Headline
posterior: \(\hat\delta\) with CrI; CrI excluding zero is a
quantitative confirmation of the McMurray–Mitchell argument that
acceleration is necessary to produce empirical vocabulary growth.

The natural ablation is \texttt{pin\_delta=0} (M2 vs.\ M1 in the
nested family, §\ref{sec:nested}), which forces the model back to
the pure accumulator. If \(\delta\) is doing real work, this ablation
should produce a substantially worse fit and force the other
parameters into pathological regions.

\subsection*{RQ2 — does frequency reconstruct word difficulty?}

The Rasch baseline (M1 in §\ref{sec:nested}) sets the coefficient on
\(\log p_j\) to 1 (the pure-accumulator assumption: each token of
word \(j\) contributes one log-token of evidence). Posterior
\(\hat\psi_j^{(s)}\) is then residual word difficulty after
unit-coefficient frequency. We report:

\begin{itemize}[leftmargin=1.3em]
\item Posterior correlation \(r(\hat\psi_j, \log p_j)\) and \(R^2\) —
  what fraction of word-level difficulty is explainable by frequency.
\item Posterior class means \(\hat\mu_c\) — does class structure
  persist after frequency is accounted for?
\item Posterior \(\beta_c\) in M4 (per-class slope on \(\log p_j\)) —
  do classes differ in how strongly frequency contributes?
\end{itemize}

The class-specific slope variant M4 is the principal RQ2 test.
\(\beta_c = 1\) corresponds to the unit-accumulator baseline;
\(\beta_c < 1\) means frequency contributes less per log-unit;
\(\beta_c \approx 0\) means frequency is uninformative once per-word
\(\psi_j\) is free.

\subsection*{RQ3 — input- vs.\ efficiency-driven variance}
\label{sec:identifiability}

The linear predictor contains \(\log r_i + \log \alpha_i\) as a sum,
unidentifiable from Wordbank alone. We break the degeneracy by pinning
\(\sigma_r^2\) at its externally-measured value (from Sperry, Hart \&
Risley, and Weisleder \& Fernald pooled). The child-level latent
\(\xi_i = \log r_i + \log \alpha_i\) is then identified from the data,
with implied distribution
\(\xi_i \sim \Norm(\mu_r,\, \sigma_r^2 + \sigma_\alpha^2)\).
Subtracting the pinned \(\sigma_r^2\) from the inferred \(\Var(\xi_i)\)
identifies \(\sigma_\alpha^2\), giving

\begin{equation}
\pi_\alpha \;\coloneqq\; \frac{\sigma_\alpha^2}{\sigma_r^2 + \sigma_\alpha^2}
\end{equation}

as a well-defined derived posterior quantity: the proportion of
child-level outcome variance attributable to learning efficiency rather
than input rate. Its complement, \(1 - \pi_\alpha\), is the input-driven
fraction.

\paragraph{What \(\pi_\alpha\) does and does not decompose.}
The full between-child variance in \(\theta_{it}\) at age \(t\) is
\[
\Var_i(\theta_{it})
= \sigma_\xi^2
+ \log_{age}^2\,\sigma_\zeta^2
+ 2\,\rho_{\xi\zeta}\,\sigma_\xi\,\sigma_\zeta\,\log_{age}.
\]
At \(t = a_0\) (the pivot, where \(\log_{age} = 0\)), this collapses
to \(\sigma_\xi^2\) and \(\pi_\alpha\) is exactly the efficiency
share. Off-pivot, the per-child slope variance \(\sigma_\zeta^2\)
contributes increasingly. With our fitted \(\sigma_\zeta \approx 3.5\)
on English, at \(\log_{age} = -1\) (e.g., 7~mo when \(a_0 = 19\)) the
slope component contributes \(\sim 5\!\times\) the intercept variance.

We do not decompose \(\sigma_\zeta\) into input-rate-growth and
efficiency-growth components, because the model treats \(\log r_i\)
as age-invariant (see ``Age-invariance of \(\log r_i\)'' in §
Assumptions). \(\zeta_i\) therefore acts uniformly on
\((\log r_i + \log \alpha_i)\) — without longitudinal observations of
\(r_i(t)\) within child, the data cannot separate it. We report
\(\pi_\alpha\) as the intercept-share at the pivot age, and treat
\(\sigma_\zeta\) as a third axis of between-child variation that
exists alongside the input/efficiency intercept decomposition.

The \texttt{drop\_slopes} ablation (\(\sigma_\zeta = 0\)) tests
whether the intercept-level \(\pi_\alpha\) is robust: forced into a
single-locus model, the data still gives \(\pi_\alpha \approx 0.92\)
on English (vs.\ \(0.91\) with slopes free), so the headline number
is robust to slope-vs-no-slope choices.

\textbf{Why the collapsed parameterization matters practically.} Sampling
\(\log r_i\) and \(\log \alpha_i\) individually — each with its own
non-centered \(\Norm(0, 1)\) prior and a scale parameter — leaves the
posterior with a flat ridge: any rotation of the pair that preserves
\(\xi_i\) has the same likelihood. Only the priors break the degeneracy,
at the cost of 2I parameters where I are identified. In HMC this
manifests as chains getting stuck in bad regions (we observed one chain
running 28+ minutes on a sim the others cleared in 3). Sampling
\(\xi_i\) directly, with \(\sigma_\alpha\) entering only through
\(\sigma_{\xi_i}^2 = \sigma_r^2 + \sigma_\alpha^2\), eliminates this
geometry entirely. Parameter recovery on a 250-child \(\times\) 150-word
synthetic dataset: \(\pi_\alpha\) posterior 0.62 [0.50, 0.72] against
true 0.61.

\textbf{Sensitivity analysis.} Refit with
\(\sigma_r^2 \in \{0.5, 1, 2\} \cdot \sigma_r^{2,\text{obs}}\).
Stability of \(\pi_\alpha\) across these indicates how much the claim
depends on trusting the external estimate.

\section{Why 2PL: one-token-one-drop is too strong}
\label{sec:twopl}

The Rasch version of (\ref{eq:eta}) (i.e.\ with \(\lambda_j \equiv 1\))
implies every word has the same logit-per-log-token gain. For a fixed
word, the predicted slope of \(\mathrm{logit}\,p_{ij}\) with respect to
age is \((1+\delta)/(a_i - s)\) — on the order of \(0.08\) logits per
month at typical observed ages. Real CDI words transition from
\(\sim\)5\% to \(\sim\)95\% acquired over roughly six months (about 6
logits), which is \(\sim\)12\(\times\) faster than the Rasch slope can
produce.

When we fit Rasch to Wordbank (500 children \(\times\) 200 items), the
posterior finds this mismatch in a predictable way: it cranks every
available steepening lever — \(\delta\) rises to \(\sim 2.3\), \(s\)
pushes to \(\sim 12\) months, and crucially \(\sigma_\alpha\) inflates
to \(\sim 2\) — implying children differ by \(\sim 50\times\) in
learning efficiency, and \(\pi_\alpha \approx 0.93\). This pattern is
stable across the s/\(\delta\) identifiability, suggesting misspecification
rather than genuine individual-difference variance.

\paragraph{The 2PL fix.} Adding word-specific discrimination
\(\lambda_j\) gives each word its own slope around threshold: the
per-word logit derivative becomes \(\lambda_j(1+\delta)/(a_i - s)\).
Words whose referents are highly diagnostic of their meaning (``ball''
when you see a ball; concrete nouns generally) can have large
\(\lambda_j\) and transition sharply; words whose meaning is more
distributional or context-dependent (function words, abstracts) have
small \(\lambda_j\) and gradual transitions. This is the standard
two-parameter logistic IRT extension and is strongly preferred over
1PL for CDI data in prior work (Kachergis, Marchman, Dale, Mankewitz,
\& Frank, 2022, \emph{JSLHR}).

\paragraph{Three accumulator interpretations of \(\lambda_j\) (all
observationally equivalent from CDI alone).}
\begin{enumerate}[leftmargin=1.3em]
\item \textbf{Information per token.} \(\lambda_j\) is the per-token
shift in the child's certainty about word \(j\). High \(\lambda_j\):
each exposure is diagnostic.
\item \textbf{Sharpness of the cognitive threshold.} \(\lambda_j\) is
the inverse width of the band over which evidence must accumulate for
acquisition to occur.
\item \textbf{ELI-to-token ratio (bridge to Mollica \& Piantadosi).} If
only a fraction of heard tokens are effective learning instances and the
rest are background noise, \(\lambda_j\) is proportional to that
fraction for word \(j\).
\end{enumerate}

\paragraph{Falsification checks.} The hierarchical prior on
\(\log \lambda_j\) (pooled via \(\sigma_\lambda\)) regularizes: one
unusual word can't take over. Two further checks protect against
freely flexible fit-to-noise:
\begin{enumerate}[leftmargin=1.3em]
\item \(\sigma_\alpha\) should \emph{decrease substantially} under 2PL.
If the Rasch inflation was misspecification dressed as person variance,
2PL should absorb it into \(\lambda_j\) and drop \(\sigma_\alpha\) toward
a more reasonable value (e.g.\ \(\sim\)0.7), sharpening our answer to
RQ4.
\item \(\lambda_j\) should correlate with measurable word properties:
higher for concrete nouns, lower for function words and abstract
adjectives. If \(\lambda_j\) is uncorrelated with any external
measurement, we have over-fit rather than learned.
\end{enumerate}

\paragraph{What this changes for the RQs.}
RQ1 (frequency-vs-difficulty) unchanged: we still report the posterior
correlation of \(\psi_j\) with \(\log p_j\). RQ2 (start time), RQ3
(\(\delta\)) unchanged at the structural level, but the extreme
posteriors observed under Rasch should moderate. RQ4 (variance
decomposition) sharpens substantially: \(\pi_\alpha\) becomes the
proportion of child variance after item-level discrimination has been
accounted for, which is the quantity we really want.

\section{Assumptions, listed}

\begin{enumerate}[leftmargin=1.3em]
\item \textbf{Pure accumulation.} Exposure accumulates linearly at a
  per-token rate independent of prior knowledge. No leveraged learning
  (Mitchell \& McMurray 2009 showed leverage cannot \emph{create}
  acceleration that is not already present).
\item \textbf{Per-word evidence gain \(\lambda_j\).} In 2PL mode each
  token contributes \(\lambda_j\) units of logit-scale evidence;
  \(\lambda_j = 1\) reduces to the Rasch ``one token, one drop'' rule.
  Per-child gain is \emph{not} modeled: all children receive the same
  \(\lambda_j\) for a given word. This is relaxable by adding a random
  slope, which we defer.
\item \textbf{Logit link with fixed scale.} Canonical Rasch/2PL IRT.
  Equivalent (up to a \(\sim\!1.7\) scale factor) to a probit model;
  approximates the Gamma cumulative of Mollica \& Piantadosi for typical
  parameter values.
\item \textbf{Cross-sectional data.} Each child observed once.
\item \textbf{CDS only.} Overheard speech ignored in v1.
\item \textbf{Same input distribution as reference studies.} Our CDI
  sample has the Sperry/HR/WF distribution of $r_i$; SES moderators not
  modeled in v1.
\item \textbf{Age-invariance of \(\log r_i\).} The model treats each
  child's input rate as a single per-kid constant, not a function of
  age. The per-child age-acceleration \((1 + \delta + \zeta_i)\) acts
  on the combined product \(r_i \cdot \alpha_i\), so any age-varying
  child-side dynamics (efficiency growth, parents talking more to
  older kids) are absorbed into \(\delta + \zeta_i\) without
  attribution. Empirically supported by within-kid regressions of
  \(\log r_{iv}^{\mathrm{obs}}\) on age in BabyView (\(N=20\) kids,
  5688 videos: mean within-kid slope \(-0.006\) logits/mo,
  median \(-0.002\)) and SEEDLingS (\(N=44\) kids, 525 LENA
  recordings: mean \(-0.011\) logits/mo). Both are at least an order
  of magnitude smaller than the \((1+\delta)\,\log_{age}\) effect over
  the same age range, so \(\sigma_\zeta\) is unlikely to be primarily
  absorbing input-rate growth. See \texttt{model/scripts/input\_rate\_vs\_age.R}.
\item \textbf{Global $s$.} One start time for all words. Per-class or
  per-word $s$ is a planned extension.
\item \textbf{CHILDES frequencies as truth.} Sampling noise in $p_j$
  ignored; defensible for common words.
\item \textbf{Lexical class only via thresholds.} Class enters through
  the hierarchical prior on $\psi_j$; class-specific frequency slopes or
  start times are variants, not baseline.
\end{enumerate}

\section{Sampling strategy: children and items}
\label{sec:sampling}

The data we fit is large enough that we always work on a subsample.
The strategy is the same across pipelines (Wordbank longitudinal,
BabyView, future Seedlings) so results are directly comparable.

\paragraph{Children.}
\begin{enumerate}[leftmargin=1.3em]
\item Filter to children with \(\geq 2\) admins (longitudinal coverage).
\item Compute each child's \emph{median admin age} as their
  representative age.
\item Bin children into \(N_{\text{age}}=4\) age quantile bins.
\item Within each bin, take the top kids by number of admins, then by
  age span. (Tie-breaking on richer longitudinal data.)
\item Top up uniformly at random if any bin is short (rare).
\end{enumerate}

This replaces an earlier, biased "top-by-priority" rule that just
maximized \(n_{\text{admin}} + \text{span}/5\); that rule overweighted
data-rich kids in narrow age windows. The age-stratified rule
guarantees coverage across the whole 16--30 mo developmental window.

\paragraph{Items.}
\begin{enumerate}[leftmargin=1.3em]
\item Filter to items with non-zero CHILDES probability \(p_j\).
\item For each lexical class, bin items into \(N_{\text{diff}}=4\)
  difficulty quartiles based on \(\log p_j\).
\item Sample uniformly within each (class, difficulty-quartile) cell
  --- 16 cells for 4 classes \(\times\) 4 quartiles.
\item Top up at random if any cell is short.
\end{enumerate}

The earlier item-sampling rule stratified only on lexical class, which
oversampled high-frequency words (those are dense in the CHILDES
distribution). The class-by-difficulty grid ensures we have hard,
medium, and easy words within \emph{each} class, which sharpens the
identification of \(\psi_j\) and \(\beta_c\) when they enter as
parameters.

\paragraph{Pipeline-shared item pool.} All Wordbank-based pipelines
(cross-sectional and longitudinal, English and Norwegian) draw from
the \emph{union of WG and WS forms}, mirroring the BabyView treatment.
Each (child, age, form) is its own admin; a child who took both WG and
WS at different ages contributes both. This extends English age coverage
from the WS-only 16--30 mo to 8--30 mo and matches the cross-pipeline
data definition.

\paragraph{Per-pipeline application.}
\begin{itemize}[leftmargin=1.3em]
\item \textbf{Wordbank longitudinal (English / Norwegian).} 200
children \(\times\) 200 items per dataset, drawn via the stratified
sampling above. English: median 3 admins, 14 mo span; Norwegian:
median 8 admins, 14 mo span (Norwegian has the much denser
Kristoffersen contributor).
\item \textbf{Wordbank cross-sectional (English).} 500 children,
treated as the most recent (child, form) admin from \texttt{long\_items.rds}.
Stratified the same way as the longitudinal pipeline.
\item \textbf{BabyView.} 20 longitudinal English subjects (the entire
available pool) \(\times\) 200 items, with input rate observed via
5\,688 videos. No child-side subsampling: we use everyone.
\end{itemize}

\section{What we deliver}

\begin{enumerate}[leftmargin=1.3em]
\item \textbf{RQ-addressed posteriors}: \(\delta\) (RQ1),
  \(\beta_c\) and \(r(\hat\psi_j, \log p_j)\) (RQ2), and
  \(\pi_\alpha\) (RQ3) — all with credible intervals.
\item \textbf{Per-word \(\psi_j\)}: absolute log-token thresholds for all
  680 CDI items.
\item \textbf{Per-child \(\log \alpha_i\)}: child-level efficiency
  estimates.
\item \textbf{Counterfactual simulator}: given hypothetical
  \((r, \alpha, a)\), predicted vocabulary size with CrI.
\item \textbf{Planned model comparisons}: LOO-CV across the baseline
  plus variants for each RQ (see §8).
\end{enumerate}

\section{Nested model family}
\label{sec:nested}

We build the model as a strictly nested progression from the simplest
possible per-word IRT (M0) up to the full-feature variant. Each step
adds one named component and answers one specific question; comparing
adjacent models (LOO-CV / WAIC, posterior shifts on the structural
parameters) attributes the marginal value of that component.

The model classes \(c \in \{1, \ldots, C\}\) of word \(j\) (lexical
category: nouns / predicates / function words / other) enter only
through hierarchical priors on \(\psi_j\) (class-mean \(\mu_c\),
class-SD \(\tau_c\)) and, in M4, through class-specific frequency
slopes \(\beta_c\). Throughout, \(\xi_i = \log r_i + \log \alpha_i\)
is the per-child intercept with prior
\(\xi_i \sim \Norm(\mu_r, \sigma_r^2 + \sigma_\alpha^2)\) and
\(\sigma_r\) externally pinned.

\paragraph{M0 — minimal IRT (no time, no frequency).}
\[
\eta_{ijt} = \xi_i + \log H - \psi_j.
\]
Per-word \(\psi_j\) absorbs all between-item variance; per-child
\(\xi_i\) absorbs all between-child intercept variance.  No age
effect, no frequency. This is the absolute null: vocabulary
trajectories are flat in age, and word difficulty is a free
per-word parameter only.

\paragraph{M1 — McMurray null (add unit-coefficient frequency and
unit-rate accumulation).}
\[
\eta_{ijt} = \xi_i + \log\!\tfrac{t-s}{a_0} + \log H + \log p_j - \psi_j.
\]
Adds the structural accumulator content: frequency enters with unit
coefficient (one log-token of evidence per token of word \(j\)
heard), and the time term grows at unit rate per log-age. This is
McMurray (2007)'s pure linear accumulator. The slope on
\(\log_{age}\) is fixed at 1, with no per-child deviation.
\textbf{M1 vs.\ M0 quantifies whether frequency and time contribute
beyond per-word \(\psi_j\) and per-child \(\xi_i\)} (RQ2 partial answer).

\paragraph{M2 — add population acceleration \(\delta\) (RQ1).}
\[
\eta_{ijt} = \xi_i + (1 + \delta)\,\log\!\tfrac{t-s}{a_0} + \log H + \log p_j - \psi_j,
\quad \delta \sim \Norm(0,\,0.5^2).
\]
Frees the \(\log_{age}\) slope as a deviation \(\delta\) above the
pure-accumulator baseline. \(\delta = 0\) recovers M1.
\textbf{M2 vs.\ M1 is the RQ1 test}: does the data require
super-linear-in-age token accumulation?

\paragraph{M3 — add per-child slope deviations \(\zeta_i\).}
\[
\eta_{ijt} = \xi_i + (1 + \delta + \zeta_i)\,\log\!\tfrac{t-s}{a_0} + \log H + \log p_j - \psi_j.
\]
Each child has their own deviation \(\zeta_i\) from the population
slope. Joint per-child prior:
\((\xi_i, \zeta_i) \sim \mathrm{MVN}(0,\Sigma)\) with LKJ on
correlation; sum-to-zero centering separates \(\delta\) from
\(\mathrm{mean}(\zeta_i)\). Requires longitudinal data for
\(\sigma_\zeta\) to be identified. \textbf{M3 vs.\ M2 measures
between-child variation in growth rate.}

\paragraph{M4 — add class-specific frequency slope \(\beta_c\) (RQ2).}
\[
\eta_{ijt} = \xi_i + (1 + \delta + \zeta_i)\,\log\!\tfrac{t-s}{a_0} + \log H + \beta_{c(j)}\,\log p_j - \psi_j,
\quad \beta_c \sim \Norm(1,\,0.5^2).
\]
Frequency now enters with class-specific weight. \(\beta_c = 1\) is
M3 (unit accumulator); \(\beta_c = 0\) means frequency is uninformative
for class \(c\) once per-word \(\psi_j\) is free.
\textbf{M4 vs.\ M3 is the RQ2 test}: do classes differ in how
strongly frequency contributes?

\paragraph{M5 — add per-word discrimination \(\lambda_j\) (2PL).}
\[
\eta_{ijt} = \lambda_j\!\left[\xi_i + (1 + \delta + \zeta_i)\,\log\!\tfrac{t-s}{a_0} + \log H + \beta_{c(j)}\,\log p_j - \psi_j\right],
\quad \log \lambda_j \sim \Norm(0,\,\sigma_\lambda^2).
\]
Per-word multiplicative scaling on the \((\theta - \beta)\) gap.
\(\sigma_\lambda \to 0\) recovers M4 (Rasch).  \textbf{M5 vs.\ M4
measures item-level sharpness variance.}

\paragraph{Robustness variants (off the main spine).}

\begin{center}
\small
\begin{tabular}{@{}lll@{}}
\toprule
Variant & Compared to & What it tests \\
\midrule
\texttt{lmm}     & M3 & Linear-in-age (\((\beta_{age} + \zeta_i)(t - a_0)\)) instead of log-linear; tests functional-form sensitivity. \\
\texttt{no\_class} & M3 & Drops class hierarchy on \(\psi_j\); tests whether class structure adds anything. \\
\texttt{free\_s}  & M3 & Frees \(s\); tests \(s\)/\(\delta\) trade-off (see Appendix~\ref{sec:s-non-id}). \\
\texttt{no\_freq} & M2 & Drops \(\log p_j\); tests whether frequency adds beyond per-word \(\psi_j\). \\
\bottomrule
\end{tabular}
\end{center}

\paragraph{The variant grammar implements the family.}
\texttt{variant\_hyperpriors()} in \texttt{model/R/helpers.R} maps a
variant name onto the prior overrides that select a specific model in
the family. Tight priors near zero pin a parameter off (e.g.,
\(\sigma_\zeta \sim \Norm(0, 0.001)\) pins \(\zeta_i \equiv 0\)); loose
priors free it. The running log of fits is in
\texttt{notes/experiments.md}.

\section{Datasets and observation channels}
\label{sec:channels}

The variant grammar (§\ref{sec:nested}) toggles \emph{model
components}. Independent of that, each dataset comes with its own
\emph{observation channels} that determine which auxiliary parameters
are appropriate. We separate the two levers in code so they can't be
mixed by accident.

\paragraph{Lever~1 — variant grammar (per-fit toggles).}
Tight-prior switches on \(\sigma_\lambda\), \(\sigma_\zeta\),
\(s\), \(\delta\). Implemented in
\texttt{variant\_hyperpriors()} (\texttt{model/R/helpers.R}); identical
across pipelines (the \texttt{long\_}, \texttt{io\_} prefixes are
stripped inside).

\paragraph{Lever~2 — dataset-level design choices (per-bundle constants).}
Properties inherent to the observation channel of a given dataset.
Baked into \texttt{stan\_data} by the \texttt{prepare\_*.R} script;
not a choice a researcher makes per fit. Two examples currently:

\begin{itemize}[leftmargin=1.3em]
\item \emph{Reactivity inflation \(\beta_{\text{react}}\).} Per-child
  observed input rate \(\log r^{\mathrm{obs}}_{iv}\) for video~\(v\) of
  child~\(i\) is
  \(\log r_i + \beta_{\text{react}} + \varepsilon_{iv}\). The bias
  \(\beta_{\text{react}}\) absorbs systematic differences between
  observed input and ``true'' input.
  \begin{itemize}
  \item \textbf{BabyView} (head-mounted video, on-camera observer
    effect): \(\beta_{\text{react}} \sim \mathcal{N}(0.4,\,0.4)\) ---
    active, with the prior centered around \(\sim\!1.5\)x inflation.
  \item \textbf{SEEDLingS} (LENA all-day audio, no on-camera observer):
    \(\beta_{\text{react}} \sim \mathcal{N}(0,\,0.001)\) --- pinned at
    zero. Passive-recording channel, so no inflation.
  \end{itemize}
\item \emph{Per-recording weighting.} Currently uniform; planned
  hook in \texttt{stan\_data\$log\_r\_obs\_weight} for downweighting
  short or contaminated recordings. Off by default.
\end{itemize}

\paragraph{Stan model files form a partial order.} Each file
\emph{adds} components on top of the previous, all toggleable via
tight priors:

\begin{center}
\small
\begin{tabular}{@{}lll@{}}
\toprule
File & Adds & Used by \\
\midrule
\texttt{log\_irt.stan} & cross-sectional baseline \(\xi_i\) & Wordbank cross-sectional fits \\
\texttt{log\_irt\_long.stan} & admin\(\to\)child indexing, \(\zeta_i\), bivariate prior & Wordbank longitudinal fits \\
\texttt{log\_irt\_io.stan} & per-recording \(\log r^{\mathrm{obs}}\) measurement layer, \(\beta_{\text{react}}\) & BabyView, SEEDLingS \\
\texttt{log\_irt\_proc.stan} \emph{(planned)} & noisy-\(\alpha\) observation from LWL RT/accuracy & Peekbank-linked Stanford datasets \\
\bottomrule
\end{tabular}
\end{center}

The result is a small partial order of model files, not a
combinatorial explosion. Adding a new dataset usually means writing
a \texttt{prepare\_<dataset>.R} script that reuses an existing Stan
file with appropriate dataset-level constants; only when a genuinely
new \emph{observation channel} appears do we add a Stan file.

\section{Comparison to prior accumulator models}
\label{sec:compare}

Five prior models in this lineage all sit, with one or two
parameterization changes, inside our nested family
(§\ref{sec:nested}). We work through each in chronological order,
showing the mapping in our notation.

\subsection*{McMurray (2007) — pure deterministic accumulator}

Each word \(j\) has a time-to-acquisition threshold \(\tau_j\) drawn
from a population distribution \(g(\tau)\). At each unit time step
every word accumulates one point, and word \(j\) is learned at the
first \(t\) with \(t \ge \tau_j\). The growth curve is the CDF of
\(g\):
\[
L_{\mathrm{McM}}(t) = \int_0^t g(\tau)\,d\tau = G(t),
\quad
L''(t) > 0 \iff g'(t) > 0.
\]
Acceleration arises only if \(g(\tau)\) puts more mass at
moderately-hard ages than at easy ones. Frequency enters
\emph{implicitly} through whatever shapes \(g(\tau)\); McMurray uses
Gaussian by CLT.

\paragraph{Mapping.} McMurray's deterministic accumulator is the
zero-noise limit of our \textbf{M1} (unit-rate accumulator with no
acceleration parameter). His \(\tau_j\) corresponds to our
log-threshold \(\psi_j\) on the time axis: in M1, the deterministic
acquisition age for word \(j\) at population mean child is
\[
t_j^{\mathrm{M1}} = a_0 \cdot \exp(\psi_j - \log p_j - \log H - \mu_r),
\]
which is the unique \(t\) where \(\theta_{\mathrm{pop}}(t) = \beta_j\).
The distribution of \(\{\psi_j\}\) across words plays the role of
\(g(\tau)\). Differences from McMurray:
\begin{enumerate}[leftmargin=1.3em]
\item \emph{Stochastic observations} — we observe Bernoulli outcomes
\(y_{ij} \sim \Bern(\sigma(\theta - \beta))\); McMurray observes
deterministic acquisition status. Our logit link is the smooth
analogue of his step function.
\item \emph{Per-child variation} — McMurray has none.
\item \emph{Frequency in the linear predictor} — explicit \(\log p_j\)
enters the linear predictor; McMurray bundles it into \(g(\tau)\).
\end{enumerate}

\subsection*{Mitchell \& McMurray (2009) — accumulator with feedback}

Continuous-time generalization with a \emph{leverage} feedback term
\(c\): the effective time argument \(a(t)\) advances faster as the
child's vocabulary \(L(t)\) grows, so
\[
a(t) = t + c\,L(t),
\quad L(t) = G(t + c\,L(t)),
\quad
\frac{dL}{dt} = \frac{g(t + cL)}{1 - c\,g(t+cL)}.
\]
Their key result: leverage shifts \(L(t)\) earlier (or later, with
\(c<0\)) but \emph{cannot create acceleration} that isn't already
present in \(g\). Acceleration still requires \(g'(\tau) > 0\) over
the relevant range. They also tested two frequency-to-difficulty
mappings (additive and multiplicative); under Zipf with exponent~1,
multiplicative gives uniform difficulties, hence linear growth.

\paragraph{Mapping.} The feedback structure
\(a(t) = t + c\,L(t)\) does not have a clean analogue in our
parameterization: it is intrinsically nonlinear in \(L\). What we
have instead is a \emph{population-level} acceleration term
\((1+\delta)\,\log_{age}\) (M2) and \emph{per-child} acceleration
deviation \(\zeta_i\) (M3), neither of which involves the child's
own vocabulary state. Mitchell \& McMurray's main \emph{conclusion}
— that acceleration must be in the model \emph{somehow} — is
operationalized in our family as the M2-vs-M1 comparison
(RQ1). We confirm their conclusion empirically with
\(\hat\delta \gg 0\) and CrI excluding zero.

\subsection*{Hidaka (2013) — closed-form AoA distributions}

Hidaka derives closed-form age-of-acquisition distributions for three
nested learning processes:
\begin{center}
\begin{tabular}{@{}lll@{}}
\toprule
Process & Dynamics & AoA distribution \\
\midrule
Cumulative & \(N\)-accumulator, constant rate \(\delta\) &
  \(\Gamma(N, \delta)\) \\
Rate-change & 1-accumulator, rate \(\delta\,t^D\) &
  Weibull \\
Cumulative + rate-change & both & Weibull--Gamma \\
\bottomrule
\end{tabular}
\end{center}
Closed-class words best fit Weibull (rate-change); content words best
fit Gamma (cumulative). \(N\) correlates with frequency;
change-of-rate exponent \(D\) correlates with imageability.

\paragraph{Mapping.} Hidaka's three families correspond to three
points in our nested family:
\begin{itemize}[leftmargin=1.3em]
\item \emph{Cumulative} = our \textbf{M1}. Constant log-age slope; AoA
  for word \(j\) at population mean child has a one-point degenerate
  distribution. With per-child variance, AoA across kids is
  approximately log-normal.
\item \emph{Rate-change} = our \textbf{M2}. Slope on \(\log_{age}\) is
  \(1 + \delta\); on the linear time scale, effective tokens grow as
  \(t^{1+\delta}\). Hidaka's exponent \(D\) corresponds (after
  reparameterization) to our \(\delta\).
\item \emph{Cumulative + rate-change} = our \textbf{M3}. Per-child
  intercept \(\xi_i\) and per-child slope deviation \(\zeta_i\) jointly
  give a per-child power-law trajectory.
\end{itemize}
The functional families differ — Hidaka's AoA is Gamma/Weibull,
ours is implicitly the convolution of a logistic CDF (item-level)
with a normal (child-level) — but both live in the same
``right-skewed positive AoA distribution'' family. Our model gains
explicit per-child structure (Hidaka fits population-averaged
\((N, \delta)\)) and embedding in a Bayesian framework with real
units.

\subsection*{Mollica \& Piantadosi (2017) — Bayesian Gamma waiting time}

Per-word generative model: accumulation begins at start age \(s_w\);
``effective learning instances'' (ELIs) arrive as a Poisson process
with rate \(\lambda_w\); word is acquired after \(k_w\) ELIs.
Acquisition age is therefore
\[
\mathrm{AoA}_w \sim s_w + \Gamma(k_w, \lambda_w),
\qquad
F_w(a) = \int_0^{a-s_w}
\frac{t^{k_w-1}\,e^{-t\lambda_w}\,\lambda_w^{k_w}}{\Gamma(k_w)}\,dt.
\]
Aggregated data \(x_{aw} \sim \mathrm{Bin}(F_w(a), N_a)\) is fit per-word
with JAGS across 13 languages. They report \(k \approx 10\),
\(\lambda \approx 0.5\)/mo, \(s \approx 2\) mo for comprehension.

\paragraph{Mapping.} Mollica's per-word \((k_w, \lambda_w)\) translates
directly to our \(\psi_j\) and the population-level rate, but with
two structural differences:
\begin{enumerate}[leftmargin=1.3em]
\item \emph{Their model is per-word; ours is per-(word, child)}.
Mollica has no \(r_i\) or \(\alpha_i\). The total rate
\(\lambda_w\) bundles input rate, child efficiency, and
frequency-of-word-\(w\) into one number. We separate
\(\lambda \approx r_i \cdot \alpha_i \cdot p_j \cdot H\), with
\(r_i\) externally pinned, enabling RQ3.
\item \emph{Gamma waiting time vs.\ logistic acquisition curve}. With
large \(k\), the Gamma cumulative converges to a normal CDF (CLT);
the logistic \(\sigma(\theta - \beta)\) is its
\(\sim\!1.7\)-scale-factor approximation. So with \(k \gtrsim 5\)
our Rasch model is approximately Mollica's Gamma waiting time, and
\(\psi_j\) maps to \(\log(k_w / \lambda_w)\) (the Gamma mean on the
log scale).
\end{enumerate}
Their joint \((k, \lambda)\) ridge is the same identifiability problem
as our \((s, \delta)\) ridge: both have one structural quantity
(the mean / median acquisition time) and a near-orthogonal shape
parameter that the data weakly constrains. Our resolution
(externally pinning \(\sigma_r\) and per-class \(\psi\) priors)
breaks the ridge in a different place than theirs.

\subsection*{Kachergis, Marchman, \& Frank (2021) — IRT framing}

Reframes the family as a 1-PL Rasch IRT model:
\[
P(y_{ij} = 1 \mid \theta_i, d_j)
= \frac{1}{1 + \exp(\theta_i + d_j)}.
\]
Fit with \texttt{lme4::glmer} on Wordbank. Item-level covariates:
expected tokens/month \((p_j \cdot H \cdot a)\), lexical class.
Person-level covariate: age. Best-fit model \texttt{m3}: lexical
class \(\times\) tok\_per\_mo + age.

\paragraph{Mapping.} Kachergis 2021 is essentially our M3 / M4 fit,
but in standardized logit units with three structural simplifications:
\begin{enumerate}[leftmargin=1.3em]
\item \emph{No external pin on \(\sigma_r\)}. Their \(\theta_i\)
absorbs both input rate and efficiency without any way to decompose
them. RQ3 is unanswerable in their formulation.
\item \emph{Age as additive person-covariate, not structural
\(\log_{age}\)}. They write \(\theta_i = \theta_{i,0} + \beta a\); we
write \(\theta_{it} = \xi_i + (1 + \delta)\,\log\!((t-s)/a_0)\).
The two are linearizations of each other for moderate age ranges,
but ours preserves units (a slope of 1 on \(\log_{age}\) is one
log-token of evidence per log-age unit).
\item \emph{No acceleration parameter}. Their model can't pose RQ1
quantitatively, although their finding that adding age to
\texttt{m\_tokmo} improves fit (i.e., age matters beyond expected
tokens) is qualitatively the same conclusion.
\end{enumerate}
Our nested family completes the standard-model program of Kachergis
2021 by (a) putting it in real token units rather than standardized
logit units, (b) adding the structural \(\delta\) for RQ1, (c)
adding the external \(\sigma_r\) pin for RQ3, and (d) extending to
input-uptake (BabyView, SEEDLingS) and processing (Peekbank-Stanford)
datasets via the observation-channel architecture.

\subsection*{Summary table}

\begin{center}
\small
\begin{tabular}{@{}lcccccc@{}}
\toprule
Model & Per-child & Per-word & Frequency & Acceleration & Bayesian & Real units \\
\midrule
McMurray 2007            & --      & \(g(\tau)\)         & implicit          & via \(g'\) only & --     & --     \\
Mitchell \& McMurray 2009 & --      & \(g(\tau)\) + \(c\)  & implicit          & via \(g'\) only & --     & --     \\
Hidaka 2013              & --      & \((N, \delta, D)\)  & via \(N\)         & via \(D\)       & --     & --     \\
Mollica \& Piantadosi 2017 & --      & \((k_w, \lambda_w)\) & implicit (in \(\lambda_w\)) & --     & \(\checkmark\) & partial \\
Kachergis et al.\ 2021   & \(\theta_i\) & \(d_j\), class    & explicit, unit-coef & --      & --     & --     \\
\textbf{Ours (M5)}       & \(\xi_i, \zeta_i\), \(\pi_\alpha\) & \(\psi_j, \lambda_j\), class & \(\beta_c\) free & \(\delta\) free & \(\checkmark\) & \(\checkmark\) \\
\bottomrule
\end{tabular}
\end{center}

\appendix

\section{Start time \(s\) and the \(s\)/\(\delta\) ridge}
\label{sec:s-non-id}

A natural fourth question — \emph{when} does accumulation start? — is
deferred to this appendix because the parameter \(s\) is
near-non-identifiable in practice. Mollica \& Piantadosi (2017)
estimated \(s \approx 2\) months for comprehension; we attempted a
similar estimate for production but found the data does not support
it, for the reason described below.

\paragraph{Identifiability in principle vs.\ practice.}
When \(s\) is freed jointly with \(\delta\), the two parameters trade
off along a near-perfect ridge: posterior
\(r(s, \delta) \approx -0.96\) on both English and Norwegian
longitudinal fits. In principle the pair is identifiable: setting
\((1+\delta)\log(t-s) = (1+\delta')\log(t-s')\) for all \(t\) and
differentiating in \(t\) gives
\((1+\delta)/(t-s) = (1+\delta')/(t-s')\), forcing
\(\delta = \delta'\) and \(s = s'\). Practically the curve-shape
difference between two points along the ridge is small. For the
English age range (16--30~mo), the ridge endpoints
\((s, \delta) = (0.5, 9.4)\) and \((3, 8)\) produce
\((1+\delta)\log\!((t-s)/a_0)\) curves that differ by only \(\sim 0.1\)
logits in shape (after subtracting a constant that \(\xi_i\)
absorbs). With Bernoulli observations and finite \(N\), that shape
difference is at the noise floor. So \(s\) is \emph{near-non-identifiable
in practice}, and the reported \(s\) posterior in any
\texttt{free\_s} variant is only a coordinate choice along the ridge
— it does not deliver an interpretable estimate of ``onset of
accumulation.'' We pin \(s \approx 0\) in M1--M5 and treat
\texttt{free\_s} as a robustness check, not a first-class estimate.

\paragraph{Reference age \(a_0\) is dataset-specific.}
\(a_0\) is the pivot at which \(\log_{age} = 0\) and so
\(\theta_{it} = \xi_i\) (per-child intercept). Setting a single
global \(a_0 = 20\) across datasets makes \(\xi_i\) language-relative
in a misleading way: when admin ages skew far from \(a_0\), the
per-kid likelihood ridge in \((\xi_i, \zeta_i)\) is tilted by an
amount proportional to \(\log(\bar t_i / a_0)\), pulling the marginal
\(r(\xi, \zeta)\) across kids in one direction or the other. We
observed this as a spurious sign flip in \(\rho_{\xi\zeta}\) between
English (median admin age 19; \(r = +0.39\)) and Norwegian (median
admin age 26; \(r = -0.32\)) on the same model. Re-evaluating each
kid's \(\xi\) at their personal mid-observation age makes the flip
vanish (Norwegian \(r\) becomes \(+0.08\)). We therefore set \(a_0\)
per dataset to the median admin age, in each \texttt{prepare\_*.R}
script. This is a reparameterization, not a model change: the
likelihood, \(\delta\), and \(\zeta_i\) are invariant; only the
\emph{interpretation} of \(\xi_i\) changes (from ``ability at 20~mo''
to ``ability at the dataset's median admin age'').

\end{document}
